\documentclass[11pt]{article}

\usepackage[a4paper,margin=1in]{geometry}
\usepackage{amsmath,amssymb,amsthm,mathtools}
\usepackage[T1]{fontenc}
\usepackage{lmodern}
\usepackage{microtype}
\usepackage{hyperref}
\usepackage{bm}

\newtheorem{theorem}{Theorem}[section]
\newtheorem{lemma}[theorem]{Lemma}
\newtheorem{proposition}[theorem]{Proposition}
\newtheorem{corollary}[theorem]{Corollary}
\newtheorem{conjecture}[theorem]{Conjecture}
\newtheorem{remark}[theorem]{Remark}
\newtheorem{definition}[theorem]{Definition}

\title{The Connectedness Threshold Problem\\
for the Clean Hatano--Nelson Family with Open Boundary Conditions}
\author{}
\date{}

\begin{document}
\maketitle

\section{Setup}

We work with the real nonnormal tridiagonal Toeplitz matrices
\[
\mathbf{A}_n=\operatorname{tridiag}(a,0,b)
=\begin{bmatrix}
0 & b \\
a & 0 & \ddots \\
& \ddots & \ddots & \ddots \\
&& \ddots & 0 & b \\
&&& a & 0
\end{bmatrix}
\in\mathbb{R}^{n\times n},
\quad n>1, \quad 0<a<b.
\]

\begin{lemma}\label{lem:real-spectrum}
Let
\[
\mathbf{D}_n:=\operatorname{diag}\bigl(1,q,q^2,\dots,q^{n-1}\bigr),
\qquad q:=\sqrt{\frac{a}{b}}.
\]
Then
\[
\mathbf{D}_n^{-1}\mathbf{A}_n\mathbf{D}_n
=\sqrt{ab}\,\operatorname{tridiag}(1,0,1).
\]
Consequently,
\[
\sigma(\mathbf{A}_n)
=\left\{2\sqrt{ab}\cos\frac{j\pi}{n+1}:j=1,\dots,n\right\}
\subset \mathbb{R}.
\]
\end{lemma}

\begin{proof}
A direct computation gives
\[
(\mathbf{D}_n^{-1}\mathbf{A}_n\mathbf{D}_n)_{j,j+1}=bq=\sqrt{ab},
\qquad
(\mathbf{D}_n^{-1}\mathbf{A}_n\mathbf{D}_n)_{j+1,j}=aq^{-1}=\sqrt{ab}.
\]
Thus \(\mathbf{D}_n^{-1}\mathbf{A}_n\mathbf{D}_n\) is the symmetric Toeplitz tridiagonal matrix
\(\sqrt{ab}\,\operatorname{tridiag}(1,0,1)\), whose spectrum is classical.
\end{proof}

For \(\varepsilon>0\), the \(\varepsilon\)-pseudospectrum of \(\mathbf{A}_n\) is
\[
\sigma_\varepsilon(\mathbf{A}_n)\coloneqq\{z\in\mathbb{C}:\|(z\mathbf{I}-\mathbf{A}_n)^{-1}\|>\varepsilon^{-1}\},
\]
where \(\|\cdot\|\) is the standard \(2\)-norm, with the convention that
\(\|(z\mathbf{I}-\mathbf{A}_n)^{-1}\|=+\infty\) for \(z\in\sigma(\mathbf{A}_n)\).
Equivalently,
\[
\sigma_\varepsilon(\mathbf{A}_n)
=\{z\in\mathbb{C}:s_{\min}(z\mathbf{I}-\mathbf{A}_n)<\varepsilon\},
\]
where \(s_{\min}(\cdot)\) denotes the smallest singular value.

\section{An Algebraic Boundary Polynomial}

For \(z=x+iy\), define
\[
\mathbf{H}_{n,\varepsilon}(x,y)
\coloneqq
\bigl((z\mathbf{I}-\mathbf{A}_n)^*(z\mathbf{I}-\mathbf{A}_n)-\varepsilon^2\mathbf{I}\bigr),
\]
and set
\[
F_{n,\varepsilon}(x,y)\coloneqq \det \mathbf{H}_{n,\varepsilon}(x,y).
\]
We shall call \(F_{n,\varepsilon}\) the \emph{\(\varepsilon\)-pseudospectral boundary polynomial}
of \(\mathbf{A}_n\).

\begin{proposition}\label{prop:boundary-polynomial}
For every \(n>1\) and \(\varepsilon>0\), the polynomial \(F_{n,\varepsilon}(x,y)\) belongs to
\(\mathbb{R}[x^2,y^2]\) and has total degree \(2n\). Moreover,
\[
\begin{aligned}
\bigl\{(x,y)\in\mathbb{R}^2:x+iy\in \partial\sigma_\varepsilon(\mathbf{A}_n)\bigr\}
&=\bigl\{(x,y)\in\mathbb{R}^2:F_{n,\varepsilon}(x,y)=0,\ \mathbf{H}_{n,\varepsilon}(x,y)\succeq 0\bigr\} \\
&\subseteq \bigl\{(x,y)\in\mathbb{R}^2:F_{n,\varepsilon}(x,y)=0\bigr\}.
\end{aligned}
\]
\end{proposition}

\begin{proof}
We first prove that \(F_{n,\varepsilon}\in\mathbb{R}[x^2,y^2]\) and that
\(\deg F_{n,\varepsilon}=2n\).

Since the entries of \(\mathbf{H}_{n,\varepsilon}(x,y)\) are polynomial in \((x,y)\),
the determinant \(F_{n,\varepsilon}(x,y)\) is a polynomial in \(\mathbb{C}[x,y]\).
For real \((x,y)\), the matrix \(\mathbf{H}_{n,\varepsilon}(x,y)\) is Hermitian, so
\(F_{n,\varepsilon}(x,y)\in\mathbb{R}\) for all \((x,y)\in\mathbb{R}^2\).
Hence \(F_{n,\varepsilon}\in\mathbb{R}[x,y]\).

Because \(\mathbf{A}_n\) is real, replacing \(y\) by \(-y\) conjugates every entry of
\(\mathbf{H}_{n,\varepsilon}(x,y)\), so
\[
\mathbf{H}_{n,\varepsilon}(x,-y)=\overline{\mathbf{H}_{n,\varepsilon}(x,y)}.
\]
Therefore
\[
F_{n,\varepsilon}(x,-y)=\overline{F_{n,\varepsilon}(x,y)}=F_{n,\varepsilon}(x,y),
\]
since \(F_{n,\varepsilon}(x,y)\) is real-valued on \(\mathbb{R}^2\).

Next, if
\[
\mathbf{J}_n\coloneqq \operatorname{diag}(1,-1,1,-1,\dots,(-1)^{n-1}),
\]
then \(\mathbf{J}_n\mathbf{A}_n\mathbf{J}_n^{-1}=-\mathbf{A}_n\), and hence
\[
\mathbf{J}_n\mathbf{H}_{n,\varepsilon}(x,y)\mathbf{J}_n^{-1}
=\mathbf{H}_{n,\varepsilon}(-x,-y).
\]
Taking determinants gives
\[
F_{n,\varepsilon}(-x,-y)=F_{n,\varepsilon}(x,y).
\]
Combining this with the evenness in \(y\) yields
\[
F_{n,\varepsilon}(-x,y)=F_{n,\varepsilon}(x,y).
\]
Thus \(F_{n,\varepsilon}\) is even in both \(x\) and \(y\), that is,
\(F_{n,\varepsilon}\in\mathbb{R}[x^2,y^2]\).

For the degree, each diagonal entry of \(\mathbf{H}_{n,\varepsilon}(x,y)\) has degree \(2\) in \((x,y)\),
all first off-diagonal entries have degree \(1\), and all second off-diagonal entries have degree \(0\).
In the determinant expansion, the identity permutation contributes the product of the diagonal entries,
whose top-degree part is
\[
(x^2+y^2)^n.
\]
Any nonidentity permutation moves at least two indices, hence uses at least two off-diagonal entries.
Replacing two diagonal entries of degree \(2\) by off-diagonal entries of degree at most \(1\)
lowers the total degree to at most \(2n-2\).
Therefore the degree-\(2n\) term comes only from the identity permutation, so
\[
\deg F_{n,\varepsilon}=2n,
\]
and in particular \(F_{n,\varepsilon}\) is not the zero polynomial.

Now let
\[
f(z)\coloneqq s_{\min}(z\mathbf{I}-\mathbf{A}_n).
\]
Since \(f\) is continuous and
\[
\sigma_\varepsilon(\mathbf{A}_n)=\{z\in\mathbb{C}:f(z)<\varepsilon\},
\]
we have the general inclusion
\[
\partial\sigma_\varepsilon(\mathbf{A}_n)\subseteq \{z\in\mathbb{C}:f(z)=\varepsilon\}.
\]

For the reverse inclusion, let \(z_0\in\mathbb{C}\) satisfy
\[
f(z_0)=\varepsilon.
\]
Then \(z_0\notin \sigma(\mathbf{A}_n)\), since \(f(z_0)=\varepsilon>0\).
Suppose, for contradiction, that
\[
z_0\notin \partial\sigma_\varepsilon(\mathbf{A}_n).
\]
Because \(f(z_0)=\varepsilon\), we have \(z_0\notin \sigma_\varepsilon(\mathbf{A}_n)\).
Since \(\sigma_\varepsilon(\mathbf{A}_n)\) is open, this means that \(z_0\) lies in the interior of its complement.
Hence there exists \(r>0\) such that
\[
B(z_0,r)\cap \sigma_\varepsilon(\mathbf{A}_n)=\varnothing.
\]
By shrinking \(r\) if necessary, we may also assume that
\[
B(z_0,r)\subseteq \rho(\mathbf{A}_n),
\]
where \(\rho(\mathbf{A}_n)\) is the resolvent set.

Define
\[
\nu(z)\coloneqq \|(z\mathbf{I}-\mathbf{A}_n)^{-1}\|,
\qquad z\in \rho(\mathbf{A}_n).
\]
On \(B(z_0,r)\) we then have
\[
\nu(z)\le \varepsilon^{-1},
\]
while at \(z_0\),
\[
\nu(z_0)=\varepsilon^{-1}.
\]
Thus \(\nu\) attains a local maximum at \(z_0\).
Since \(z\mapsto (z\mathbf{I}-\mathbf{A}_n)^{-1}\) is analytic on \(\rho(\mathbf{A}_n)\),
the scalar function \(\nu\) is subharmonic there. By the maximum principle for subharmonic functions,
\(\nu\) must be constant on \(B(z_0,r)\). Consequently,
\[
f(z)=\nu(z)^{-1}=\varepsilon
\qquad \text{for all } z\in B(z_0,r).
\]
Hence
\[
F_{n,\varepsilon}(x,y)=0
\qquad \text{for all } (x,y)\ \text{with}\ x+iy\in B(z_0,r),
\]
because one singular value of \(z\mathbf{I}-\mathbf{A}_n\) equals \(\varepsilon\) throughout that open set.
This contradicts the fact that \(F_{n,\varepsilon}\) is a nonzero polynomial.
Therefore
\[
\{z\in\mathbb{C}:f(z)=\varepsilon\}\subseteq \partial\sigma_\varepsilon(\mathbf{A}_n),
\]
and so
\[
\partial\sigma_\varepsilon(\mathbf{A}_n)=\{z\in\mathbb{C}:f(z)=\varepsilon\}.
\]

Finally, the eigenvalues of
\[
\mathbf{H}_{n,\varepsilon}(x,y)
=\bigl((z\mathbf{I}-\mathbf{A}_n)^*(z\mathbf{I}-\mathbf{A}_n)-\varepsilon^2\mathbf{I}\bigr)
\]
are precisely \(s_j(z\mathbf{I}-\mathbf{A}_n)^2-\varepsilon^2\), where \(s_j\) are the singular values
of \(z\mathbf{I}-\mathbf{A}_n\). Hence
\[
s_{\min}(z\mathbf{I}-\mathbf{A}_n)=\varepsilon
\iff
\lambda_{\min}\bigl(\mathbf{H}_{n,\varepsilon}(x,y)\bigr)=0
\iff
\bigl(F_{n,\varepsilon}(x,y)=0\ \text{and}\ \mathbf{H}_{n,\varepsilon}(x,y)\succeq 0\bigr).
\]
This proves the stated boundary characterization.
\end{proof}

\begin{proposition}[No isolated boundary points]\label{prop:no-isolated-boundary-points}
For every \(n>1\) and \(\varepsilon>0\), the boundary \(\partial\sigma_\varepsilon(\mathbf{A}_n)\) has no isolated points.
\end{proposition}

\begin{proof}
Set \(\Omega:=\sigma_\varepsilon(\mathbf{A}_n)\). Then \(\Omega\) is a nonempty proper open subset of \(\mathbb{C}\).
Suppose that \(z_0\in\partial\Omega\) is isolated in \(\partial\Omega\). Then there exists \(r>0\) such that
\[
\partial\Omega\cap B(z_0,r)=\{z_0\}.
\]
Hence the punctured disk \(B(z_0,r)\setminus\{z_0\}\) is the disjoint union of
\[
\Omega\cap (B(z_0,r)\setminus\{z_0\})
\quad\text{and}\quad
\Omega^c\cap (B(z_0,r)\setminus\{z_0\}).
\]
Both sets are nonempty because \(z_0\in\partial\Omega\), and both are open in the relative topology of the punctured disk because there are no boundary points other than \(z_0\) in \(B(z_0,r)\). This contradicts the connectedness of \(B(z_0,r)\setminus\{z_0\}\).
\end{proof}

\begin{remark}\label{rem:not-exact-algebraic-curve}
In general, the full zero set
\[
\bigl\{(x,y)\in\mathbb{R}^2:F_{n,\varepsilon}(x,y)=0\bigr\}
\]
is larger than the true boundary
\[
\bigl\{(x,y)\in\mathbb{R}^2:x+iy\in \partial\sigma_\varepsilon(\mathbf{A}_n)\bigr\}.
\]
The equation \(F_{n,\varepsilon}(x,y)=0\) only says that \emph{some} singular value of
\(z\mathbf{I}-\mathbf{A}_n\) equals \(\varepsilon\); the boundary is the branch on which the
\emph{smallest} singular value equals \(\varepsilon\).
\end{remark}

\begin{proposition}[Discarded algebraic branches lie inside the pseudospectrum]\label{prop:discarded-branches-inside}
If \((x,y)\in\mathbb{R}^2\) satisfies
\[
F_{n,\varepsilon}(x,y)=0
\qquad\text{and}\qquad
\mathbf{H}_{n,\varepsilon}(x,y)\not\succeq 0,
\]
then
\[
x+iy\in \sigma_\varepsilon(\mathbf{A}_n).
\]
Equivalently,
\[
\bigl\{(x,y):F_{n,\varepsilon}(x,y)=0\bigr\}
\setminus
\bigl\{(x,y):x+iy\in \partial\sigma_\varepsilon(\mathbf{A}_n)\bigr\}
\subseteq
\bigl\{(x,y):x+iy\in \sigma_\varepsilon(\mathbf{A}_n)\bigr\}.
\]
\end{proposition}

\begin{proof}
Let \(z=x+iy\). Since \(F_{n,\varepsilon}(x,y)=0\), the Hermitian matrix \(\mathbf{H}_{n,\varepsilon}(x,y)\) has a zero eigenvalue. Since it is not positive semidefinite, it also has a negative eigenvalue. Thus
\[
\lambda_{\min}\bigl(\mathbf{H}_{n,\varepsilon}(x,y)\bigr)<0.
\]
By definition of \(\mathbf{H}_{n,\varepsilon}(x,y)\),
\[
\lambda_{\min}\bigl(\mathbf{H}_{n,\varepsilon}(x,y)\bigr)
=s_{\min}(z\mathbf{I}-\mathbf{A}_n)^2-\varepsilon^2.
\]
Hence \(s_{\min}(z\mathbf{I}-\mathbf{A}_n)<\varepsilon\), that is, \(z\in\sigma_\varepsilon(\mathbf{A}_n)\).
\end{proof}

\begin{proposition}\label{prop:explicit-H}
Set
\[
\tau\coloneqq x^2+y^2-\varepsilon^2,
\qquad
c\coloneqq az+b\bar z=(a+b)x-i(b-a)y,
\qquad
p\coloneqq ab,
\qquad
d\coloneqq \tau+a^2+b^2.
\]
Then
\[
\mathbf{H}_{n,\varepsilon}(x,y)
=(|z|^2-\varepsilon^2)\mathbf{I}-\bar z\,\mathbf{A}_n-z\,\mathbf{A}_n^{\mathsf{T}}+\mathbf{A}_n^{\mathsf{T}}\mathbf{A}_n,
\]
and its entries are given by
\[
\bigl(\mathbf{H}_{n,\varepsilon}(x,y)\bigr)_{jk}
=
\begin{cases}
\tau+a^2, & j=k=1,\\
\tau+b^2, & j=k=n,\\
d, & j=k\in\{2,\dots,n-1\},\\
-c, & k=j+1,\\
-\bar c, & j=k+1,\\
p, & |j-k|=2,\\
0, & |j-k|>2.
\end{cases}
\]
Equivalently, for \(n\ge 3\),
\[
\mathbf{H}_{n,\varepsilon}(x,y)=
\begin{bmatrix}
\tau+a^2 & -c & p \\
-\bar c & d & -c & \ddots \\
p & -\bar c & d & \ddots & \ddots \\
& \ddots & \ddots & \ddots & \ddots & \ddots \\
&& \ddots & \ddots & d & -c & p \\
&&& \ddots & -\bar c & d & -c \\
&&&& p & -\bar c & \tau+b^2
\end{bmatrix},
\]
with the obvious truncation when \(n=2\), where
\[
\mathbf{H}_{2,\varepsilon}(x,y)=
\begin{bmatrix}
\tau+a^2 & -c \\
-\bar c & \tau+b^2
\end{bmatrix}.
\]
\end{proposition}

\begin{proof}
Expanding the product gives
\[
\mathbf{H}_{n,\varepsilon}(x,y)
=(\bar z\mathbf{I}-\mathbf{A}_n^{\mathsf{T}})(z\mathbf{I}-\mathbf{A}_n)-\varepsilon^2\mathbf{I}
=(|z|^2-\varepsilon^2)\mathbf{I}-\bar z\,\mathbf{A}_n-z\,\mathbf{A}_n^{\mathsf{T}}+\mathbf{A}_n^{\mathsf{T}}\mathbf{A}_n.
\]
Now \(-\bar z\,\mathbf{A}_n-z\,\mathbf{A}_n^{\mathsf{T}}\) contributes the first off-diagonals
\(-\bigl(az+b\bar z\bigr)=-c\) and \(-\bigl(a\bar z+bz\bigr)=-\bar c\).
Moreover, \(\mathbf{A}_n^{\mathsf{T}}\mathbf{A}_n\) has diagonal entries
\(a^2,a^2+b^2,\dots,a^2+b^2,b^2\), has second off-diagonals equal to \(ab=p\),
and has all first off-diagonal entries equal to \(0\). The stated formula follows.
\end{proof}

\begin{corollary}\label{cor:Toeplitz-reduction}
Let \(\Delta_0\coloneqq 1\), and for \(m\ge 1\) let
\[
\Delta_m\coloneqq \det \mathbf{T}_m(d,c,p),
\]
where
\[
\mathbf{T}_1(d,c,p)\coloneqq [d],
\qquad
\mathbf{T}_2(d,c,p)\coloneqq
\begin{bmatrix}
d & -c \\
-\bar c & d
\end{bmatrix},
\]
and for \(m\ge 3\),
\[
\mathbf{T}_m(d,c,p)
=\begin{bmatrix}
d & -c & p \\
-\bar c & d & -c & \ddots \\
p & -\bar c & d & \ddots & \ddots \\
& \ddots & \ddots & \ddots & \ddots & \ddots \\
&& \ddots & \ddots & d & -c & p \\
&&& \ddots & -\bar c & d & -c \\
&&&& p & -\bar c & d
\end{bmatrix}.
\]
Then, for every \(n\ge 2\),
\[
F_{n,\varepsilon}(x,y)
=\Delta_n-(a^2+b^2)\Delta_{n-1}+a^2b^2\Delta_{n-2}.
\]
\end{corollary}

\begin{proof}
By Proposition~\ref{prop:explicit-H},
\[
\mathbf{H}_{n,\varepsilon}(x,y)=\mathbf{T}_n(d,c,p)-b^2\mathbf{E}_{11}-a^2\mathbf{E}_{nn},
\]
where \(\mathbf{E}_{jj}\) denotes the matrix with a single \(1\) in the \((j,j)\) entry.
Since the determinant is affine in each of the two modified diagonal entries, we obtain
\[
\det \mathbf{H}_{n,\varepsilon}
=\det \mathbf{T}_n
-b^2\det \mathbf{T}_n^{(1)}
-a^2\det \mathbf{T}_n^{(n)}
+a^2b^2\det \mathbf{T}_n^{(1,n)},
\]
where \(\mathbf{T}_n^{(1)}\) and \(\mathbf{T}_n^{(n)}\) are obtained from \(\mathbf{T}_n\) by deleting,
respectively, the first row and column and the last row and column, and
\(\mathbf{T}_n^{(1,n)}\) is obtained by deleting both the first and the last row and column.
Because \(\mathbf{T}_n\) is Toeplitz, these minors are precisely
\(\mathbf{T}_{n-1}(d,c,p)\), \(\mathbf{T}_{n-1}(d,c,p)\), and \(\mathbf{T}_{n-2}(d,c,p)\)
(with the convention \(\Delta_0=1\) when \(n=2\)).
Therefore
\[
F_{n,\varepsilon}(x,y)
=\Delta_n-(a^2+b^2)\Delta_{n-1}+a^2b^2\Delta_{n-2}.
\]
\end{proof}

\begin{remark}\label{rem:quartic-closed-form}
The Toeplitz determinants \(\Delta_m\) admit a closed form in \(m\). Let
\(\rho_1,\rho_2,\rho_3,\rho_4\) be the roots of the quartic equation
\[
p\rho^4-c\rho^3+d\rho^2-\bar c\rho+p=0,
\]
that is,
\[
ab\rho^4-\bigl((a+b)x-i(b-a)y\bigr)\rho^3
+\bigl(x^2+y^2-\varepsilon^2+a^2+b^2\bigr)\rho^2
-\bigl((a+b)x+i(b-a)y\bigr)\rho
+ab=0.
\]
When the roots are distinct, one convenient closed form is
\[
\Delta_m
=p^m
\frac{
\det
\begin{bmatrix}
1&1&1&1\\
\rho_1&\rho_2&\rho_3&\rho_4\\
\rho_1^{m+2}&\rho_2^{m+2}&\rho_3^{m+2}&\rho_4^{m+2}\\
\rho_1^{m+3}&\rho_2^{m+3}&\rho_3^{m+3}&\rho_4^{m+3}
\end{bmatrix}
}{
\det
\begin{bmatrix}
1&1&1&1\\
\rho_1&\rho_2&\rho_3&\rho_4\\
\rho_1^2&\rho_2^2&\rho_3^2&\rho_4^2\\
\rho_1^3&\rho_2^3&\rho_3^3&\rho_4^3
\end{bmatrix}
},
\]
with the repeated-root case obtained by the usual confluent limit. Hence
\[
F_{n,\varepsilon}(x,y)
=\Delta_n-(a^2+b^2)\Delta_{n-1}+a^2b^2\Delta_{n-2}
\]
provides a closed form in \(n\).
\end{remark}

\section{Reduction of Simply Connectedness to a Vertical-Slice Problem}

\begin{theorem}[Vertical-slice criterion for simple connectedness]\label{thm:vertical-slice-criterion}
Fix \(n>1\) and \(\varepsilon>0\). Assume that for every connected component \(\Omega\) of
\(\sigma_\varepsilon(\mathbf{A}_n)\) and every \(x\in\mathbb{R}\), the set
\[
\{y\ge 0:x+iy\in \partial\Omega\}
\]
has at most one element. Then every connected component of \(\sigma_\varepsilon(\mathbf{A}_n)\) is simply connected.
\end{theorem}

\begin{proof}
Let \(\Omega\) be a connected component of \(\sigma_\varepsilon(\mathbf{A}_n)\).
Since \(\|(z\mathbf{I}-\mathbf{A}_n)^{-1}\|\to 0\) as \(|z|\to\infty\), the set
\(\sigma_\varepsilon(\mathbf{A}_n)\), hence also \(\Omega\), is bounded.

We first claim that \(\Omega\) contains an eigenvalue of \(\mathbf{A}_n\). Suppose not.
Then \(\overline\Omega\cap\sigma(\mathbf{A}_n)=\varnothing\), because every eigenvalue belongs to
\(\sigma_\varepsilon(\mathbf{A}_n)\) and therefore lies in the interior of some connected component.
Thus the function
\[
\nu(z):=\|(z\mathbf{I}-\mathbf{A}_n)^{-1}\|
\]
is continuous on \(\overline\Omega\) and subharmonic on a neighborhood of \(\overline\Omega\).
Since \(\Omega\) is a connected component of the open set \(\sigma_\varepsilon(\mathbf{A}_n)\), we have
\(\partial\Omega\subseteq\partial\sigma_\varepsilon(\mathbf{A}_n)\), so Proposition~\ref{prop:boundary-polynomial} gives
\[
\nu(z)=\varepsilon^{-1}
\qquad \text{for all } z\in\partial\Omega.
\]
On the other hand,
\[
\nu(z)>\varepsilon^{-1}
\qquad \text{for all } z\in\Omega.
\]
This contradicts the maximum principle for subharmonic functions. Hence \(\Omega\) contains an eigenvalue of \(\mathbf{A}_n\).
By Lemma~\ref{lem:real-spectrum}, every eigenvalue of \(\mathbf{A}_n\) is real, so
\[
\Omega\cap\mathbb{R}\neq\varnothing.
\]

Because \(\mathbf{A}_n\) is real, \(\sigma_\varepsilon(\mathbf{A}_n)\) is invariant under complex conjugation. Set
\[
\Omega^{\sharp}:=\{\bar z:z\in\Omega\}.
\]
Then \(\Omega^{\sharp}\) is also a connected component of \(\sigma_\varepsilon(\mathbf{A}_n)\). Since
\(\Omega\cap\mathbb{R}\neq\varnothing\), the components \(\Omega\) and \(\Omega^{\sharp}\) intersect, and therefore
\[
\Omega=\Omega^{\sharp}.
\]
Hence each vertical fiber
\[
\Omega_x:=\{y\in\mathbb{R}:x+iy\in\Omega\}
\]
is symmetric with respect to \(y\mapsto -y\).

Fix \(x\in\mathbb{R}\). The set \(\Omega_x\) is an open subset of \(\mathbb{R}\). By hypothesis, its nonnegative boundary set
\[
\{y\ge 0:x+iy\in\partial\Omega\}
\]
has cardinality at most one. Since \(\Omega_x\) is open, bounded, and symmetric, it follows that either
\(\Omega_x=\varnothing\) or else there exists \(h(x)>0\) such that
\[
\Omega_x=(-h(x),h(x)).
\]
Indeed, any other symmetric open subset of \(\mathbb{R}\) would contribute at least two nonnegative boundary points.

Now define
\[
I:=\{x\in\mathbb{R}:\Omega_x\neq\varnothing\}.
\]
Then \(I\) is open, and
\[
\Omega=\{x+iy:x\in I,\ |y|<h(x)\}.
\]
If \(I\) were disconnected, then \(\Omega\) would split as a union of two disjoint nonempty open subsets indexed by the two open components of \(I\), contradicting connectedness. Hence \(I\) is a connected open subset of \(\mathbb{R}\), and so \(I\) is an open interval.

Finally, the homotopy
\[
\Phi_t(x+iy):=x+i(1-t)y,
\qquad 0\le t\le 1,
\]
leaves \(\Omega\) invariant, because each nonempty fiber is the interval \((-h(x),h(x))\).
Thus \(\Phi_t\) is a deformation retraction of \(\Omega\) onto the interval \(I\). Therefore \(\Omega\) is contractible, hence simply connected.
\end{proof}

\section{Paired-Cell Reduction, Stripped Parity Weyl Maps, and Type-1 Normalization}

Retain the notation
\[
\tau:=x^2+y^2-\varepsilon^2,
\qquad
c:=az+b\bar z=(a+b)x-i(b-a)y,
\qquad
p:=ab,
\qquad
d:=\tau+a^2+b^2,
\]
with \(z=x+iy\). Define the normalized parameters
\[
\alpha:=\frac{d}{p}
=\frac{x^2+y^2-\varepsilon^2+a^2+b^2}{ab},
\qquad
\alpha_L:=\frac{\tau+a^2}{p}
=\alpha-\frac{b}{a},
\qquad
\alpha_R:=\frac{\tau+b^2}{p}
=\alpha-\frac{a}{b},
\]
and
\[
\beta:=\frac{c}{p}
=\frac{(a+b)x-i(b-a)y}{ab}.
\]
Set
\[
\mathbf{Q}:=
\begin{bmatrix}
\alpha&-\beta\\
-\bar\beta&\alpha
\end{bmatrix},
\qquad
\mathbf{Q}_L:=
\begin{bmatrix}
\alpha_L&-\beta\\
-\bar\beta&\alpha
\end{bmatrix},
\qquad
\mathbf{Q}_R:=
\begin{bmatrix}
\alpha&-\beta\\
-\bar\beta&\alpha_R
\end{bmatrix},
\]
\[
\mathbf{L}:=
\begin{bmatrix}
1&0\\
-\beta&1
\end{bmatrix},
\qquad
\mathbf{b}_\beta:=
\begin{bmatrix}
1\\
-\beta
\end{bmatrix}.
\]

\begin{proposition}[Even case: normalized block Jacobi form]\label{prop:even-normalized-block-Jacobi}
Let \(n=2m\). Then \(p^{-1}\mathbf{H}_{2m,\varepsilon}(x,y)\), viewed in the consecutive \(2\times2\) block decomposition
\[
\mathbb{C}^{2m}=\mathbb{C}^2\oplus\cdots\oplus \mathbb{C}^2,
\]
is the Hermitian block Jacobi matrix
\[
\mathbf{K}_m^{\mathrm{ev}}(x,y)=
\begin{bmatrix}
\mathbf{Q}_L & \mathbf{L} \\
\mathbf{L}^* & \mathbf{Q} & \mathbf{L} \\
& \ddots & \ddots & \ddots \\
&& \mathbf{L}^* & \mathbf{Q} & \mathbf{L} \\
&&& \mathbf{L}^* & \mathbf{Q}_R
\end{bmatrix},
\]
with the obvious truncation when \(m=1\), namely
\[
\mathbf{K}_1^{\mathrm{ev}}(x,y)=
\begin{bmatrix}
\alpha_L&-\beta\\
-\bar\beta&\alpha_R
\end{bmatrix}.
\]
\end{proposition}

\begin{proof}
Write \(\boldsymbol{\psi}=(\psi_1,\dots,\psi_{2m})^{\mathsf T}\) and group the entries into cells
\[
\boldsymbol{\eta}_k:=
\begin{bmatrix}
\psi_{2k-1}\\
\psi_{2k}
\end{bmatrix},
\qquad k=1,\dots,m.
\]
Using Proposition~\ref{prop:explicit-H}, the equations
\[
\mathbf{H}_{2m,\varepsilon}(x,y)\boldsymbol{\psi}=0
\]
become
\[
\widetilde{\mathbf{D}}_L\boldsymbol{\eta}_1+\widetilde{\mathbf{T}}\boldsymbol{\eta}_2=0,
\]
\[
\widetilde{\mathbf{T}}^*\boldsymbol{\eta}_{k-1}+\widetilde{\mathbf{D}}_0\boldsymbol{\eta}_k+\widetilde{\mathbf{T}}\boldsymbol{\eta}_{k+1}=0,
\qquad k=2,\dots,m-1,
\]
\[
\widetilde{\mathbf{T}}^*\boldsymbol{\eta}_{m-1}+\widetilde{\mathbf{D}}_R\boldsymbol{\eta}_m=0,
\]
where
\[
\widetilde{\mathbf{D}}_0=
\begin{bmatrix}
d&-c\\
-\bar c&d
\end{bmatrix},
\qquad
\widetilde{\mathbf{T}}=
\begin{bmatrix}
p&0\\
-c&p
\end{bmatrix},
\]
\[
\widetilde{\mathbf{D}}_L=
\begin{bmatrix}
\tau+a^2&-c\\
-\bar c&d
\end{bmatrix},
\qquad
\widetilde{\mathbf{D}}_R=
\begin{bmatrix}
d&-c\\
-\bar c&\tau+b^2
\end{bmatrix}.
\]
Dividing by \(p=ab>0\) gives exactly the displayed block matrix.
\end{proof}

\begin{proposition}[Odd case: normalized block Jacobi form]\label{prop:odd-normalized-block-Jacobi}
Assume \(n=2m+1\). In the decomposition
\[
\mathbb{C}^{2m+1}
=
\mathbb{C}^2\oplus\cdots\oplus\mathbb{C}^2\oplus\mathbb{C},
\]
the normalized matrix \(p^{-1}\mathbf{H}_{2m+1,\varepsilon}(x,y)\) is the Hermitian block matrix
\[
\mathbf{K}_m^{\mathrm{odd}}(x,y)=
\begin{bmatrix}
\mathbf{Q}_L & \mathbf{L} \\
\mathbf{L}^* & \mathbf{Q} & \mathbf{L} \\
& \ddots & \ddots & \ddots \\
&& \mathbf{L}^* & \mathbf{Q} & \mathbf{b}_\beta \\
&&& \mathbf{b}_\beta^* & \alpha_R
\end{bmatrix}.
\]
\end{proposition}

\begin{proof}
Write \(\boldsymbol{\psi}=(\psi_1,\dots,\psi_{2m+1})^{\mathsf T}\), set
\[
\boldsymbol{\eta}_k:=
\begin{bmatrix}
\psi_{2k-1}\\
\psi_{2k}
\end{bmatrix},
\qquad k=1,\dots,m,
\qquad
\xi:=\psi_{2m+1},
\]
and again use Proposition~\ref{prop:explicit-H}. The equations
\[
\mathbf{H}_{2m+1,\varepsilon}(x,y)\boldsymbol{\psi}=0
\]
become
\[
\widetilde{\mathbf{D}}_L\boldsymbol{\eta}_1+\widetilde{\mathbf{T}}\boldsymbol{\eta}_2=0,
\]
\[
\widetilde{\mathbf{T}}^*\boldsymbol{\eta}_{k-1}+\widetilde{\mathbf{D}}_0\boldsymbol{\eta}_k+\widetilde{\mathbf{T}}\boldsymbol{\eta}_{k+1}=0,
\qquad k=2,\dots,m-1,
\]
\[
\widetilde{\mathbf{T}}^*\boldsymbol{\eta}_{m-1}+\widetilde{\mathbf{D}}_0\boldsymbol{\eta}_m+\widetilde{\mathbf{b}}\,\xi=0,
\qquad
\widetilde{\mathbf{b}}:=
\begin{bmatrix}
p\\
-c
\end{bmatrix},
\]
and
\[
\widetilde{\mathbf{b}}^{\,*}\boldsymbol{\eta}_m+(\tau+b^2)\xi=0.
\]
Dividing by \(p\) gives the claimed form, with \(\mathbf{b}_\beta=\widetilde{\mathbf{b}}/p\).
\end{proof}

\begin{lemma}[Positivity of intermediate Schur complements]\label{lem:intermediate-Schur-positive}
Let
\[
\mathbf{M}=
\begin{bmatrix}
\mathbf{X} & \mathbf{Y}\\
\mathbf{Y}^* & \mathbf{Z}
\end{bmatrix}
\]
be Hermitian with \(\mathbf{M}\succeq 0\), where \(\mathbf{X}\) and \(\mathbf{Y}\) are square of the same size and \(\mathbf{Y}\) is invertible.
Then \(\mathbf{X}\succ 0\).
\end{lemma}

\begin{proof}
If \(\mathbf{u}\in\ker \mathbf{X}\), then
\[
0\le
\begin{bmatrix}
\mathbf{u}\\ t\mathbf{v}
\end{bmatrix}^*
\mathbf{M}
\begin{bmatrix}
\mathbf{u}\\ t\mathbf{v}
\end{bmatrix}
=2t\,\Re(\mathbf{u}^*\mathbf{Y}\mathbf{v})+t^2\mathbf{v}^*\mathbf{Z}\mathbf{v}
\]
for all \(t\in\mathbb{R}\) and all \(\mathbf{v}\). Hence \(\mathbf{u}^*\mathbf{Y}=0\), so \(\mathbf{Y}^*\mathbf{u}=0\). Since \(\mathbf{Y}\) is invertible, \(\mathbf{Y}^*\) has trivial kernel, and therefore \(\mathbf{u}=0\). Thus \(\mathbf{X}\) is positive definite.
\end{proof}

\begin{proposition}[Even stripped parity Weyl recursion]\label{prop:even-stripped-parity-Weyl}
Assume \(n=2m\ge 4\). Define recursively
\[
\mathbf{W}_1:=\mathbf{Q}_L,
\qquad
\mathbf{W}_{k+1}:=\mathbf{Q}-\mathbf{L}^*\mathbf{W}_k^{-1}\mathbf{L},
\qquad k=1,\dots,m-2,
\]
whenever the inverse exists, and define the terminal matrix
\[
\boldsymbol{\Xi}_m:=\mathbf{Q}_R-\mathbf{L}^*\mathbf{W}_{m-1}^{-1}\mathbf{L}.
\]
Then
\[
\det \mathbf{H}_{2m,\varepsilon}(x,y)=0
\iff
\det \boldsymbol{\Xi}_m(x,y)=0.
\]
Moreover,
\[
\mathbf{H}_{2m,\varepsilon}(x,y)\succeq 0
\iff
\mathbf{W}_1(x,y),\dots,\mathbf{W}_{m-1}(x,y)\succ 0
\ \text{and}\
\boldsymbol{\Xi}_m(x,y)\succeq 0.
\]
Consequently,
\[
x+iy\in \partial\sigma_\varepsilon(\mathbf{A}_{2m})
\iff
\mathbf{W}_1,\dots,\mathbf{W}_{m-1}\succ 0,
\qquad
\boldsymbol{\Xi}_m\succeq 0,
\qquad
\det \boldsymbol{\Xi}_m=0.
\]
\end{proposition}

\begin{proof}
Apply block Gaussian elimination to \(\mathbf{K}_m^{\mathrm{ev}}\). Since every off-diagonal block is the invertible matrix \(\mathbf{L}\), Lemma~\ref{lem:intermediate-Schur-positive} shows that positive semidefiniteness of \(\mathbf{K}_m^{\mathrm{ev}}\) forces each leading Schur complement to be positive definite until the final step. Thus the standard Schur-complement recursion is well defined and yields precisely \(\mathbf{W}_1,\dots,\mathbf{W}_{m-1},\boldsymbol{\Xi}_m\). This proves both the determinant identity and the semidefiniteness criterion. Since \(\mathbf{K}_m^{\mathrm{ev}}=p^{-1}\mathbf{H}_{2m,\varepsilon}(x,y)\) with \(p>0\), the same statements hold for \(\mathbf{H}_{2m,\varepsilon}(x,y)\).
\end{proof}

\begin{proposition}[Odd stripped parity Weyl recursion]\label{prop:odd-stripped-parity-Weyl}
Assume \(n=2m+1\ge 3\). Define recursively
\[
\mathbf{W}_1:=\mathbf{Q}_L,
\qquad
\mathbf{W}_{k+1}:=\mathbf{Q}-\mathbf{L}^*\mathbf{W}_k^{-1}\mathbf{L},
\qquad k=1,\dots,m-1,
\]
whenever the inverse exists, and define the terminal scalar
\[
\vartheta_m:=\alpha_R-\mathbf{b}_\beta^*\mathbf{W}_m^{-1}\mathbf{b}_\beta.
\]
Then
\[
\det \mathbf{H}_{2m+1,\varepsilon}(x,y)=0
\iff
\vartheta_m(x,y)=0.
\]
Moreover,
\[
\mathbf{H}_{2m+1,\varepsilon}(x,y)\succeq 0
\iff
\mathbf{W}_1(x,y),\dots,\mathbf{W}_m(x,y)\succ 0
\ \text{and}\
\vartheta_m(x,y)\ge 0.
\]
Consequently,
\[
x+iy\in \partial\sigma_\varepsilon(\mathbf{A}_{2m+1})
\iff
\mathbf{W}_1,\dots,\mathbf{W}_m\succ 0,
\qquad
\vartheta_m=0.
\]
\end{proposition}

\begin{proof}
Apply the same Schur-complement elimination to \(\mathbf{K}_m^{\mathrm{odd}}\). The first \(m\) steps again use the invertibility of \(\mathbf{L}\) and Lemma~\ref{lem:intermediate-Schur-positive}. The last Schur complement is the scalar \(\vartheta_m\). Multiplication by the positive scalar \(p\) does not affect determinants or semidefiniteness.
\end{proof}

\begin{proposition}[Type-1 normalization and the orbit matrix]\label{prop:type-1-normalization}
Set
\[
\mathbf{R}_1:=\mathbf{I}_2,
\qquad
\mathbf{R}_{k+1}:=\mathbf{L}^{-1}\mathbf{R}_k^{-*},
\qquad k\ge 1,
\]
and
\[
\mathbf{M}_\beta:=\mathbf{L}^{-1}\mathbf{L}^*
=
\begin{bmatrix}
1&-\bar\beta\\
\beta&1-|\beta|^2
\end{bmatrix}.
\]
Then
\[
\mathbf{R}_{k+2}=\mathbf{M}_\beta \mathbf{R}_k,
\qquad k\ge 1.
\]
Hence
\[
\mathbf{R}_{2j+1}=\mathbf{M}_\beta^j,
\qquad
\mathbf{R}_{2j}=\mathbf{M}_\beta^{j-1}\mathbf{L}^{-1},
\qquad j\ge 1.
\]
In particular,
\[
\det \mathbf{M}_\beta=1,
\qquad
\operatorname{tr}\mathbf{M}_\beta=2-|\beta|^2.
\]

Moreover, \(\mathbf{K}_m^{\mathrm{ev}}\) is congruent to the type-1 block Jacobi matrix
\[
\widetilde{\mathbf{K}}_m^{\mathrm{ev}}=
\begin{bmatrix}
\mathbf{B}_1^{(L)} & \mathbf{I}_2 \\
\mathbf{I}_2 & \mathbf{B}_2 & \mathbf{I}_2 \\
& \ddots & \ddots & \ddots \\
&& \mathbf{I}_2 & \mathbf{B}_{m-1} & \mathbf{I}_2 \\
&&& \mathbf{I}_2 & \mathbf{B}_m^{(R)}
\end{bmatrix},
\]
where
\[
\mathbf{B}_1^{(L)}=\mathbf{Q}_L,
\qquad
\mathbf{B}_k=\mathbf{R}_k^*\mathbf{Q}\mathbf{R}_k\quad (2\le k\le m-1),
\qquad
\mathbf{B}_m^{(R)}=\mathbf{R}_m^*\mathbf{Q}_R\mathbf{R}_m.
\]

Likewise, \(\mathbf{K}_m^{\mathrm{odd}}\) is congruent to
\[
\widetilde{\mathbf{K}}_m^{\mathrm{odd}}=
\begin{bmatrix}
\mathbf{B}_1^{(L)} & \mathbf{I}_2 \\
\mathbf{I}_2 & \mathbf{B}_2 & \mathbf{I}_2 \\
& \ddots & \ddots & \ddots \\
&& \mathbf{I}_2 & \mathbf{B}_m & \mathbf{u}_m \\
&&& \mathbf{u}_m^* & \alpha_R
\end{bmatrix},
\]
where
\[
\mathbf{B}_1^{(L)}=\mathbf{Q}_L,
\qquad
\mathbf{B}_k=\mathbf{R}_k^*\mathbf{Q}\mathbf{R}_k\quad (2\le k\le m),
\qquad
\mathbf{u}_m=\mathbf{R}_m^*\mathbf{b}_\beta.
\]
\end{proposition}

\begin{proof}
Since
\[
\mathbf{L}^{-1}=
\begin{bmatrix}
1&0\\
\beta&1
\end{bmatrix},
\]
we have
\[
\mathbf{R}_{k+2}
=
\mathbf{L}^{-1}\mathbf{R}_{k+1}^{-*}
=
\mathbf{L}^{-1}\bigl(\mathbf{L}^{-1}\mathbf{R}_k^{-*}\bigr)^{-*}
=
\mathbf{L}^{-1}\mathbf{L}^*\mathbf{R}_k
=
\mathbf{M}_\beta \mathbf{R}_k.
\]
The explicit formulas for \(\mathbf{R}_{2j+1}\) and \(\mathbf{R}_{2j}\) follow by induction. The determinant and trace of \(\mathbf{M}_\beta\) are immediate.

Now let
\[
\mathbf{R}_m^{\mathrm{ev}}:=\operatorname{diag}(\mathbf{R}_1,\dots,\mathbf{R}_m),
\qquad
\mathbf{R}_m^{\mathrm{odd}}:=\operatorname{diag}(\mathbf{R}_1,\dots,\mathbf{R}_m,1).
\]
By construction,
\[
\mathbf{R}_k^*\mathbf{L}\mathbf{R}_{k+1}
=
\mathbf{R}_k^*\mathbf{L}\mathbf{L}^{-1}\mathbf{R}_k^{-*}
=
\mathbf{I}_2.
\]
Therefore
\[
(\mathbf{R}_m^{\mathrm{ev}})^*\,\mathbf{K}_m^{\mathrm{ev}}\,\mathbf{R}_m^{\mathrm{ev}}
=
\widetilde{\mathbf{K}}_m^{\mathrm{ev}},
\qquad
(\mathbf{R}_m^{\mathrm{odd}})^*\,\mathbf{K}_m^{\mathrm{odd}}\,\mathbf{R}_m^{\mathrm{odd}}
=
\widetilde{\mathbf{K}}_m^{\mathrm{odd}},
\]
with the stated blocks.
\end{proof}

\begin{corollary}[Affine relation along a vertical line]\label{cor:alpha-beta-affine}
Fix \(x\in\mathbb{R}\). Then
\[
\alpha(y)=\frac{x^2+y^2-\varepsilon^2+a^2+b^2}{ab},
\qquad
\beta(y)=\frac{(a+b)x-i(b-a)y}{ab},
\]
and
\[
|\beta(y)|^2
=
\frac{(b-a)^2}{ab}\,\alpha(y)+\mu_x,
\]
where
\[
\mu_x
=
\frac{4x^2}{ab}
+
\frac{(b-a)^2(\varepsilon^2-a^2-b^2)}{a^2b^2}.
\]
Equivalently,
\[
|\beta(y)|^2
=
\frac{(a+b)^2x^2+(b-a)^2y^2}{a^2b^2}.
\]
\end{corollary}

\begin{proof}
The explicit formula for \(\beta(y)\) gives the second identity immediately. Substituting
\[
y^2=ab\,\alpha(y)-x^2+\varepsilon^2-a^2-b^2
\]
into that identity yields the affine relation.
\end{proof}

\begin{remark}[Local scalarization in the even case]\label{rem:local-scalarization-even}
Assume \(n=2m\ge 4\). On the open semialgebraic set
\[
\mathcal{D}_{m,+}^{\mathrm{ev}}
:=
\bigl\{(x,y)\in\mathbb{R}^2:
\mathbf{W}_1,\dots,\mathbf{W}_{m-1}\succ 0,
\ (\boldsymbol{\Xi}_m)_{11}>0
\bigr\},
\]
define
\[
\sigma_m(x,y)
:=
(\boldsymbol{\Xi}_m)_{22}
-
\frac{|(\boldsymbol{\Xi}_m)_{12}|^2}{(\boldsymbol{\Xi}_m)_{11}}
=
\frac{\det \boldsymbol{\Xi}_m(x,y)}{(\boldsymbol{\Xi}_m)_{11}}.
\]
Then, on \(\mathcal{D}_{m,+}^{\mathrm{ev}}\),
\[
\boldsymbol{\Xi}_m\succeq 0
\ \text{and}\
\det \boldsymbol{\Xi}_m=0
\iff
\sigma_m=0.
\]
Thus, on any region where \((\boldsymbol{\Xi}_m)_{11}>0\), the even boundary condition is equivalent to the scalar terminal equation \(\sigma_m=0\) together with the positivity constraints on \(\mathbf{W}_1,\dots,\mathbf{W}_{m-1}\).
\end{remark}

\begin{corollary}[Weyl reformulation of the vertical-slice problem]\label{cor:remaining-vertical-slice-problem}
Theorem~\ref{thm:vertical-slice-criterion} reduces the simple-connectedness problem to a vertical-slice uniqueness statement.
Using Propositions~\ref{prop:even-stripped-parity-Weyl} and \ref{prop:odd-stripped-parity-Weyl}, this uniqueness statement can be rewritten entirely in terms of the stripped parity Weyl data:

\begin{itemize}
\item if \(n=2m\), it suffices to show that for every connected component \(\Omega\) of \(\sigma_\varepsilon(\mathbf{A}_{2m})\) and every \(x\in\mathbb{R}\), the set
\[
\bigl\{y\ge 0:x+iy\in\partial\Omega,\ \mathbf{W}_1(x,y),\dots,\mathbf{W}_{m-1}(x,y)\succ 0,\ \boldsymbol{\Xi}_m(x,y)\succeq 0,\ \det \boldsymbol{\Xi}_m(x,y)=0\bigr\}
\]
has at most one element;

\item if \(n=2m+1\), it suffices to show that for every connected component \(\Omega\) of \(\sigma_\varepsilon(\mathbf{A}_{2m+1})\) and every \(x\in\mathbb{R}\), the set
\[
\bigl\{y\ge 0:x+iy\in\partial\Omega,\ \mathbf{W}_1(x,y),\dots,\mathbf{W}_m(x,y)\succ 0,\ \vartheta_m(x,y)=0\bigr\}
\]
has at most one element.
\end{itemize}

Under either of these uniqueness statements, every connected component of \(\sigma_\varepsilon(\mathbf{A}_n)\) is simply connected.
\end{corollary}

\begin{proof}
This is immediate from Theorem~\ref{thm:vertical-slice-criterion} together with the boundary characterizations in Propositions~\ref{prop:even-stripped-parity-Weyl} and \ref{prop:odd-stripped-parity-Weyl}.
\end{proof}

\end{document}